\begin{abstract}
Programming languages are commonly designed around English keywords, English
diagnostics, and English-language documentation. While many tools support
localized documentation or messages, the surface syntax of programming
languages remains overwhelmingly Anglophone. This paper presents \plat, a small
interpreted language whose keywords, diagnostics, examples, and naming
conventions are designed around Limburgian linguistic identity. Rather than
proposing new semantic mechanisms, \plat{} serves as a case study in culturally
localized language design. We discuss design tradeoffs involving dialect
selection, ASCII compatibility, learnability, authenticity, and the boundary
between educational language, esolang, and practical tool. The paper argues
that regional-language programming systems can function as instruments of
cultural visibility and digital language preservation, even when their
computational core is intentionally conventional.
\end{abstract}
